\documentclass[11pt]{article}

\usepackage[letterpaper,margin=1in]{geometry}
\usepackage[hidelinks]{hyperref}
\usepackage{titlesec}
\usepackage[dvipsnames]{xcolor}
\usepackage{needspace}
\usepackage{parskip}
\usepackage{enumitem}

\setlength{\parindent}{0pt}
\setlist[itemize]{leftmargin=1.2em,itemsep=0.2em,topsep=0.2em}
\titleformat{\section}{\LARGE\bfseries\color{RoyalBlue}}{}{0em}{}[\color{RoyalBlue}\titlerule]
\titlespacing*{\section}{0pt}{1.0em}{0.6em}

\newcommand{\eduentry}[4]{%
\begin{tabular*}{\textwidth}{@{}l@{\extracolsep{\fill}}r@{}}
\textbf{#1} & #2 \\
#3 & #4 \\
\end{tabular*}\vspace{0.6em}
}

\newcommand{\refentry}[5]{%
\textbf{#1}\\
#2\\
#3\\
#4\\
#5
}

\newcommand{\cvsection}[2][12]{%
\Needspace{#1\baselineskip}%
\section*{#2}%
}

\newenvironment{paperentry}{%
\samepage
}{%
\vspace{0.6em}%
\par
}

\begin{document}
\pagestyle{plain}

\begin{center}
{\LARGE \textbf{Paige Rowberry}}\\[0.3em]
1655 Campus Center Drive, Salt Lake City, Utah 84102 \textbar{} Office: BuC 440\\
\href{mailto:Paige.Rowberry@utah.edu}{Paige.Rowberry@utah.edu} \textbar{}
\href{https://paigenn.github.io/paigerowberry/}{paigenn.github.io/paigerowberry} \textbar{} Updated: 09/2026
\end{center}

\cvsection{Education}
\eduentry{University of Utah}{Salt Lake City, Utah}{Ph.D. in Finance}{2021--Current}
\eduentry{Brigham Young University}{Provo, Utah}{B.S. in Economics}{2021}
\eduentry{Brigham Young University}{Provo, Utah}{B.S. in Finance}{2021}

\cvsection{Research Interests}
\textbf{Corporate Finance, Firm Location, and Market Microstructure}

\cvsection{Working Papers}
\begin{paperentry}
\textbf{\href{https://papers.ssrn.com/sol3/papers.cfm?abstract_id=5258696}{Valuing Agglomeration}}\\
Job Market Paper

Abstract: State and local governments spend billions to attract corporate headquarters, promising localized economic spillovers. Using a novel dataset of U.S. headquarters relocations, I evaluate these claims by measuring how equity markets price incumbent firms when a peer's headquarters arrival is announced. Same-industry incumbents gain value, while incumbents in different industries do not show the same positive valuation response. The return premium decays with distance from the relocating firm and strengthens with product similarity. Adding a same-industry peer also increases local patenting and establishment growth. Therefore, relocation incentives are most defensible when they attract companies to locations with existing industry peers.
\end{paperentry}

\begin{paperentry}
\textbf{\href{https://papers.ssrn.com/sol3/papers.cfm?abstract_id=4834348}{Persistent Spillovers from Temporary Pandemic Restrictions}}\\
with Andra Ghent and Matthew Spiegel\\
\emph{Revise and resubmit (third round), Journal of Financial Economics}

Abstract: Pandemic restrictions targeted non-essential businesses that required in-person contact to operate. We document the impact of these restrictions county-by-county on directly regulated businesses and on businesses not subject to restrictions. Through 2023, two years after nearly all restrictions were lifted, a one-standard-deviation increase in restrictions increased establishment exits by 2\% and decreased net job creation by 13\% relative to pre-period medians. These adverse effects extended beyond regulated firms to those exempt from regulations. We attribute this transmission at least in part to the loss of face-to-face interactions.
\end{paperentry}

\clearpage

\begin{paperentry}
\textbf{\href{https://papers.ssrn.com/sol3/papers.cfm?abstract_id=5857262}{Does Market Fragmentation Improve Price Efficiency?}}\\
with Jonathan Brogaard and Anne Haubo Dyhrberg

Abstract: In equity markets the last three decades have been defined by an explosion in trading fragmentation. However, the impact of fragmentation on price efficiency is not settled. Using the collapse of FTX and coin delistings on Binance as natural experiments, we test the predictions of Chen and Duffie (2021). We find supporting evidence that market fragmentation increases price efficiency and that these gains have diminishing returns: going from two to three venues yields a larger improvement than going from nine to ten. We show that order-splitting is the mechanism that drives these results.
\end{paperentry}

\cvsection{Peer-Reviewed Publications}
\begin{paperentry}
\textbf{New Evidence on the Cycle in the Women, Infants, and Children Program: What Happens When Benefits Expire}\\
with Marianne Bitler, Jason Cook, and Seojung Oh\\
\emph{Published, National Tax Journal, 2024}

Abstract: We provide evidence of a benefit redemption cycle in the Special Supplemental Nutrition Program for Women, Infants, and Children (WIC). Using novel administrative data on item-level redemptions, we show that unlike SNAP, WIC redemptions peak at both the beginning and the end of the month. Beginning-of-the-month excess redemptions are concentrated among more fully redeemed items such as infant formula, while end-of-the-month excess redemptions are concentrated among less fully redeemed items such as infant meats. We document that a substantial share of beneficiaries go at least one month without redeeming anything and discuss how administrative burdens may drive the WIC cycle.
\end{paperentry}

\cvsection{Awards and Honors}
\begin{tabular*}{\textwidth}{@{}l@{\extracolsep{\fill}}r@{}}
Marriner S. Eccles Graduate Fellowship in Political Economy & 2025 \\
\end{tabular*}

\cvsection{Teaching Experience}
\textbf{University of Utah}\\
FINAN 4030: Corporate Finance (sole instructor)\hfill Spring 2024

\textbf{Brigham Young University}\\
Finance Research Seminar (teaching assistant)\hfill Fall 2020\\
Principles of Finance (teaching assistant)\hfill Spring 2019

\clearpage
\cvsection[12]{References}
\begin{minipage}[t]{0.48\textwidth}\raggedright
\refentry
{Andra Ghent (Chair)}
{Ivory-Boyer Chair in Real Estate}
{Department of Finance, David Eccles School of Business}
{University of Utah}
{Email: \href{mailto:Andra.Ghent@Eccles.Utah.edu}{Andra.Ghent@Eccles.Utah.edu}}

\vspace{1em}
\refentry
{Jason Cook}
{Assistant Professor}
{Marriner S. Eccles Institute, David Eccles School of Business}
{University of Utah}
{Email: \href{mailto:Jason.Cook@Eccles.Utah.edu}{Jason.Cook@Eccles.Utah.edu}}

\vspace{1em}
\refentry
{Sara Malik}
{Assistant Professor}
{School of Accounting, David Eccles School of Business}
{University of Utah}
{Email: \href{mailto:Sara.Malik@Eccles.Utah.edu}{Sara.Malik@Eccles.Utah.edu}}
\end{minipage}\hfill
\begin{minipage}[t]{0.48\textwidth}\raggedright
\refentry
{Matthew Spiegel}
{Professor of Finance}
{Director of Graduate Studies, Yale School of Management}
{Yale University}
{Email: \href{mailto:matthew.spiegel@yale.edu}{matthew.spiegel@yale.edu}}

\vspace{1em}
\refentry
{Jonathan Brogaard}
{Kendall D. Garff Chaired Professor}
{Department of Finance, David Eccles School of Business}
{University of Utah}
{Email: \href{mailto:brogaardj@eccles.utah.edu}{brogaardj@eccles.utah.edu}}
\end{minipage}

\end{document}
